\chapter*{}
\thispagestyle{empty}
\cleardoublepage

%\thispagestyle{empty}

\input{portada/portada_2}



\cleardoublepage
\thispagestyle{empty}

\begin{center}
{\large\bfseries \myTitle: \mySubTitle}\\
\end{center}

%\vspace{0.7cm}
\noindent{\textbf{Palabras clave}: Generación Aumentada por Recuperación, RAG, Grandes Modelos de Lenguaje, LLM, Modelos de Embedding, Inteligencia Artificial, IA, Ollama, OpenAI, Mistral, LangChain, Chroma, Streamlit, BOE, DANA}\\

\vspace{0.7cm}
\noindent{\textbf{Resumen}}\\

En este Trabajo de Fin de Máster, hemos diseñado, desarrollado y evaluado un sistema de generación aumentada por recuperación o RAG. Aplicando diversas técnicas en la construcción, configuración y evaluación, hemos conseguido obtener un sistema capaz de responder preguntas con bastante precisión y de forma muy eficiente.\\

\noindent La base de un sistema RAG son los grandes modelos de lenguaje o LLM. Estos están capacitados para responder preguntas generalistas, pero no funcionan de igual modo con datos con los que no han sido entrenados o privados. Para solucionar esta problemática, un sistema RAG necesita acceder a una base de conocimiento externa para construir una respuesta específica usando estos nuevos datos. Estas bases de conocimiento suelen ser bases de datos vectoriales que guardan la relación semántica entre los documentos, y los recuperan de forma rápida y eficiente ante una consulta de usuario.\\

\noindent En este proyecto hemos implementado nuestro propio sistema RAG junto a una interfaz web en la que podemos realizar consultas. También se ha elaborado un amplio conjunto de pruebas de rendimiento de las diferentes alternativas con las que podemos configurar un sistema RAG (modelos de lenguaje, embedding o estrategias de búsqueda) con el objetivo de averiguar la mejor combinación de elementos. Para ello, se han definido cuatro agentes críticos encargados de evaluar las respuestas generadas por el RAG.\\

\noindent Durante la realización de este proyecto hemos sopesado diferentes opciones para realizar nuestro sistema RAG. Hemos explorado y descartado plataformas con las que solíamos trabajar dado a que se han quedado obsoletas ante el auge de los grandes modelos de lenguaje actualmente. Librerías como LangChain (diseñada para trabajar con modelos de lenguaje), o Chroma (base de datos vectorial); han facilitado el desarrollo del proyecto gracias a su flexibilidad y adaptabilidad a diferentes modelos de lenguaje: desde modelos locales de Ollama, externos y gratuitos de Mistral o de pago con OpenAI.\\

\noindent Por último, la base de conocimiento usada para el proyecto contiene documentos del Boletín Oficial Español (BOE) con referencias a lo sucedido en la DANA (Depresión Aislada en Niveles Altos) de finales del año 2024. Cuando el proyecto comenzó, este evento estaba muy reciente y pensamos que era una buena idea la de proporcionar una plataforma de ayuda a los ciudadanos con dudas o preguntas sobre la DANA. Nuestro sistema RAG esta diseñado solo para este cometido, y a través de su interfaz web, responde con claridad y precisión cualquier duda siempre que tenga datos sólidos para ello.

\cleardoublepage


\thispagestyle{empty}


\begin{center}
{\large\bfseries \myTitleEn: \mySubTitleEn}\\
\end{center}

%\vspace{0.7cm}
\noindent{\textbf{Keywords}: Retrieval augmented generation, RAG, large language models, LLM, embedding models, artificial intelligence, AI, Ollama, OpenAI, Mistral, LangChain, Chroma, Streamlit}\\

\vspace{0.7cm}
\noindent{\textbf{Abstract}}\\

In this Master's Thesis, we have designed, developed and evaluated a retrieval augmented generation (RAG) system. By applying various techniques in construction, configuration and evaluation, we have achieved a system capable of answering question with accuracy and efficiency.\\

\noindent The foundation of an RAG system lies in Large Language Models (LLMs). These models are trained to answer general questions but may not perform equally well with unfamiliar or private data. To address this issue, a RAG system needs access to an external knowledge base to construct specific responses using new data. These knowledge bases typically consist of vector databases that store semantic relationships between documents and retrieve them quickly and efficiently in response to user queries.\\

\noindent In this project, we implemented our own RAG system alongside a web interface for conducting queries. We also develop extensive performance testing of different configurations for RAG systems (language models, embeddings, search strategies) to determine the optimal combination of elements. To facilitate this, we defined four critical agents tasked with evaluating the responses generated by the RAG system.\\

\noindent Throughout the course of this project, we considered various options for implementing our RAG system. We explored and discarded platforms that had become obsolete in light of the current dominance of large language models. Libraries such as LangChain (designed for working with language models) and Chroma (a vector database) facilitated project development due to their flexibility and adaptability to different language models: from local models like Ollama's, to external free models like Mistral's, and paid models like those from OpenAI.\\

\noindent Finally, the knowledge base used for the project contains documents from the Official State Gazette (BOE) with references to the events of the late 2024 cut-off low (DANA) weather phenomenon. At the project's inception, this event was recent, prompting us to create a platform to assist citizens with questions about DANA. Our RAG system is designed specifically for this purpose, providing clear and precise answers through its web interface, provided it has solid data to draw upon.

\chapter*{}
\thispagestyle{empty}

\noindent\rule[-1ex]{\textwidth}{2pt}\\[4.5ex]

Yo, \textbf{\myName}, alumno de la titulación \myDegree de la \textbf{\myFaculty \ de la \myUni}, con DNI 76652136J, autorizo la ubicación de la siguiente copia de mi Trabajo Fin de Máster en la biblioteca del centro para que pueda ser
consultada por las personas que lo deseen.

\vspace{6cm}

\noindent Fdo: \myName

\vspace{2cm}

\begin{flushright}
\myTime
\end{flushright}


\chapter*{}
\thispagestyle{empty}

\noindent\rule[-1ex]{\textwidth}{2pt}\\[4.5ex]

D. \textbf{\myProf}, Profesor del Departamento de Lenguaje y Sistemas Informáticos de la Universidad de Granada.

\vspace{0.5cm}

\vspace{0.5cm}

\textbf{Informan:}

\vspace{0.5cm}

Que el presente trabajo, titulado \textit{\textbf{\myTitle: \mySubTitle}}, ha sido realizado bajo su supervisión por \textbf{\myName \ }, y autorizamos la defensa de dicho trabajo ante el tribunal que corresponda.

\vspace{0.5cm}

Y para que conste, expiden y firman el presente informe en Granada a \myTime.

\vspace{1cm}

\textbf{Los directores:}

\vspace{5cm}

\noindent \textbf{\myProf}

\chapter*{Agradecimientos}
\thispagestyle{empty}

\vspace{1cm}


Quiero agradecer especialmente a \myProf\ y \myOtherProf\ por su guía, disponibilidad y valiosos consejos a lo largo de este trabajo.

\vspace{0.5cm}

\noindent A mis compañeros y compañeras de máster, gracias por el apoyo constante, las discusiones enriquecedoras y la camaradería durante este proceso. También agradezco profundamente a mi familia, por su paciencia, ánimo y confianza en mí.

\vspace{0.5cm}

\noindent Finalmente, a todas las personas que, de una forma u otra, contribuyeron a la realización de este trabajo: gracias.